\section{Veins and Integration with SUMO through TraCI}
\label{env:veins}

Veins is the framework that enables the integration between network simulation in OMNeT++ and road traffic simulation, making possible a bidirectional co-simulation in which mobility and communication influence one another. This setup follows the methodological line proposed in the literature on coupled network-traffic simulation, which shows that the analysis of IVC/V2X applications is much more credible when channel behavior and traffic behavior are simulated jointly \cite{sommer2011bidirectionally,eckhoff2012multichannel}.

In the workflow of this thesis, Veins provides the coupling infrastructure between OMNeT++ and SUMO, along with the basic components for vehicular networking scenarios. The link between the two simulators is realized through TraCI, the \emph{Traffic Control Interface}, which keeps vehicle mobility and network node state aligned during execution. OMNeT++ receives traffic information from SUMO, while communication and control logic can influence scenario behavior, making it possible to study the interaction between platoon dynamics and communication quality.

As for SUMO, its role is that of a microscopic traffic simulator, widely adopted in the ITS community for defining reproducible and configurable road scenarios \cite{alvarez2018microscopic}. In this thesis, SUMO is not used as a standalone tool, but as an integrated component of the co-simulation: its main function is to provide the mobility dynamics of the scenario, while communication dynamics and part of the application logic are handled on the OMNeT++/Veins/PLEXE side.

In summary, Veins and TraCI are the operational bridge between networking and traffic. Their presence in the development environment is what makes it possible, in the following chapters, to study the behavior of the multi-interface system in realistic scenarios and with metrics directly relevant to platooning, rather than limiting the analysis to a purely networking perspective.
